\pagestyle{fancy}\thispagestyle{plain}

\rule{6.8in}{1pt}\vskip.1in

{\large Assignment 5: \csname topic05\endcsname}\smallskip

\rule{6.8in}{1pt}\vskip.1in

\textit{Due by \HWdueTime,   eastern, on \dueDay, \csname dateWeek05\endcsname}

\rule{6.8in}{1pt}\vskip.1in

{\bf\textit{Suggested readings  for this problem set:}}
\begin{itemize}[itemsep=1pt]
\item Section 2.1, p.~49-57; 
\item Section 2.2, p.~61-65 (stop at DeMorgan's laws)
\end{itemize}

All readings are from \bookAuthors, \emph{\bookName}.

\rule{6.8in}{1pt}\vskip.1in

{\bf\textit{Extra Practice Problems}} from \emph{\bookName} (not to turn in):

\begin{itemize}[itemsep=1pt]
% Section 2.1, \#7,9,12,14,21,D1 \#13, \#18a, 19a-c, 20e-f, §2.2, \#14, \#26 (2.1) 6,7,9,10,12,13

\item With answers or hints (in the back of the book); many of these are calculations; do as many as you need to do to understand the definitions:
 \begin{enumerate}
 \item Section 2.1, \#1(adg), 2(adg), 4(adg), 5(a), 7(a), 8(ae), 9(adf), 10(a), 18(acf), 19(ad), 20(ae), 21
 \item Section 2.2, \#1(adgj), 2(ad), 4(ad), 5(ad), 7(a), 9(ad), 14(a),
\end{enumerate}

\item Without answers:
 \begin{enumerate}
 \item Section 2.1, 13, 14, 15, 16, 
 \item Section 2.2, \#1-12
 \end{enumerate}
  
\item \href{https://dmzb.github.io/teaching/2024Spring220/handouts/220-H06-sets-I.pdf}{Handout 6}
\end{itemize}

\rule{6.8in}{1pt}\vskip.1in

 {\bf\textit{Assignment:}} due \dueDay, \csname dateWeek05\endcsname, \HWdueTime, via Gradescope (\gradescopeCode):

\rule{6.8in}{1pt}\vskip.1in

\begin{enumerate}

\item Let $A = \{0,1,2\}$. Let $B = \{1,\{2\},3\}$. Which of the following statements are true? (No justification is needed.)

\begin{multicols}{3}  
 \begin{enumerate}
 \item $0 \in A$
 \item $\{0\} \subseteq A$
 \item $\{0\} \in A$      
\item $A = B$
\item $2 \in B$
\item $\{2\} \in B$
  \item $\{2\} \subseteq B$
  \item $2 \in A \cup B$
      \item $2 \in A \cap B$
      % \item $2 \in B - A$
  \item $2 \in A - B$        
   
  \end{enumerate}
\end{multicols}
  
\item Let $A = \{ n \in \mathbb{Z} \,|\, n \text{ is a multiple of } 4\}$ and $B = \{ n \in \mathbb{Z} \,|\, n^2 \text{ is a multiple of } 4\}$ 
 \begin{enumerate}
 \item Prove or disprove: $A \subseteq B$. 
 \item Prove or disprove: $B \subseteq A$. 
 \end{enumerate}

\item Recall that $(a,b) = \{x : x \in \mathbb{R} \,|\, a < x < b\}$. Prove or disprove each of the following:
 \begin{enumerate}
 \item $(-1,1) \subseteq (-2,2)$. 
 \item $(-1,2) \subseteq (-2,1)$.
 \end{enumerate}

 \item Prove that $(-10,5] \cap [0,10] = [0,5]$.
 
\item Prove that if $A \not \subseteq C$ then  $A \not \subseteq B$ or $B \not \subseteq C$. (Remember that you can reference and use lemmas and other things we proved in class.)

 \item Let $A,B$ be sets. Prove each of the following:
 \begin{enumerate}
 \item $A \cap B \subseteq A$;
 \item $A \cap \emptyset = \emptyset$;
 \end{enumerate}
  
\item Prove that $A \cup (A \cap B) = A.$

% \item Let $A = \{ n \in \mathbb{Z} | n = 8t+7 \text{ for some } t \in \mathbb{Z} \}$ and $B = \{ n \in \mathbb{Z} | n = 4t + 3 \text{ for some } t \in \mathbb{Z}\}$ 
% \begin{enumerate}
% \item Prove of disprove: $A \subseteq B$. 
% \item Prove of disprove: $B \subseteq A$. 
% \end{enumerate}
\item Let $n$ and $m$ be integers. Prove that if $m$ divides $n$ then $n\mathbb{Z} \subseteq m\mathbb{Z}$.

 

\end{enumerate}

%%% Local Variables: 
%%% mode: latex
%%% TeX-master: "../assignments-S24-220"
%%% End: 